% ==============================================================================
% RESUMEN / ABSTRACT
% ==============================================================================
% Instrucciones de la plantilla:
% - Redactar en español e inglés (exactamente 1 párrafo para cada idioma).
% - Contenido sugerido:
%   * Objetivo principal de la práctica de laboratorio.
%   * Metodología empleada y herramientas/bibliotecas utilizadas.
%   * Principales resultados alcanzados y aprendizajes obtenidos.
% ==============================================================================

\begin{abstract}
En esta práctica de laboratorio se desarrollaron e implementaron los conceptos fundamentales de la computación gráfica moderna mediante el uso del lenguaje C++ y la biblioteca gráfica OpenGL 3.3 en perfil núcleo. A lo largo de la práctica se configuraron los cauces de renderizado a través de búferes de vértices (VBO) y objetos de arreglos de vértices (VAO), implementando transformaciones afines y control interactivo en tiempo real. Los resultados obtenidos demuestran la viabilidad de la técnica empleada para generar escenas tridimensionales con alto rendimiento, identificando los retos asociados al manejo de memoria en la GPU y la sincronización de cuadros.

\vspace{0.3cm}
\noindent\textbf{Palabras clave:} Computación gráfica, OpenGL, VAO, VBO, shaders, transformaciones geométricas.
\end{abstract}

\vspace{0.5cm}

\renewcommand{\abstractname}{Abstract}
\begin{abstract}
In this laboratory session, fundamental concepts of modern computer graphics were implemented using C++ and the OpenGL 3.3 core profile. Throughout the practice, rendering pipelines were set up using vertex buffer objects (VBO) and vertex array objects (VAO), successfully applying affine transformations and real-time interactive controls. The obtained results demonstrate the effectiveness of the proposed technique for high-performance 3D scene generation, highlighting the challenges related to GPU memory management and frame synchronization.

\vspace{0.3cm}
\noindent\textbf{Keywords:} Computer graphics, OpenGL, VAO, VBO, shaders, geometric transformations.
\end{abstract}

\vspace{0.6cm}
